\chapter*{Wstęp}
\phantomsection
\addcontentsline{toc}{chapter}{Wstęp}
\markboth{WSTĘP}{WSTĘP}

\section*{Dla kogo jest ta książka}

Umiesz wywołać \code{softmax} i nie umiesz powiedzieć, dlaczego najpierw
odejmuje się od niego jego własne maksimum. Wytrenowałeś model i nigdy nie
wyprowadziłeś gradientu, za którym szedł. Czytasz artykuł, rozumiesz każde
zdanie wokół wzorów i wzory pomijasz.

To nie jest wada charakteru ani nic wyjątkowego. Tak się dzieje, gdy dziedzina
rośnie szybciej niż jej podręczniki, a matematyka, na której się opiera, jest
wykładana w kolejności zaprojektowanej dla matematyków. Ta książka jest
napisana dla inżyniera w takiej sytuacji: zawodowo sprawnego, formalnie
niedowyposażonego i świadomego tego.

Nie zakłada \textbf{niczego}. Ani algebry, ani trygonometrii, ani rachunku
różniczkowego. Część~I zaczyna się od arytmetyki i to nie jest figura
retoryczna. Jeśli masz wykształcenie ścisłe, nie szkodzi \dash{} Quizy istnieją
właśnie po to, żebyś mógł pominąć to, co już umiesz \dash{} ale podłoga jest
naprawdę podłogą, bo książka, która twierdzi, że zaczyna od zera, a potem po
cichu zakłada potęgi, zaczęła od jedynki.

\section*{Po co jest}

Dla sprawności, nie dla ścisłości. Po jej przerobieniu wyprowadzisz mechanizm
uwagi \emph{scaled dot-product}, powiesz dokładnie, co robi czynnik
$1/\sqrt{d_k}$ i co się psuje bez niego, przeczytasz funkcję straty jako
wyrażenie, a nie jako nazwę, i spojrzysz na tabelę wyników, żeby powiedzieć,
czy różnica w niej jest prawdziwa.

Nie nauczysz się dowodzić. To zamierzone ograniczenie metody i jest ono
powtórzone w każdym miejscu, w którym daje o sobie znać, w ramce nazywającej
to, co zostało przyjęte bez dowodu, i wskazującej, gdzie ten dowód znaleźć.
Książka udająca, że daje jedno i drugie w jednym tomie, nie daje żadnego.

\section*{Dlaczego format jest dziwny}

Ta książka korzysta z metody nauczania programowanego K.~A. Strouda, z
\emph{Engineering Mathematics} (1970). Materiał jest zupełnie inny; maszyneria
jest jego i została użyta dlatego, że jest jedynym szeroko stosowanym formatem
w nauczaniu matematyki, który kiedykolwiek zmierzono wobec grupy kontrolnej.

Maszyneria jest prosta. Książka to ciąg ponumerowanych \emph{ramek}. Prawie
każda żąda odpowiedzi, a odpowiedź znajduje się na początku następnej ramki,
więc masz się zadeklarować, zanim sprawdzisz. W obrębie jednego programu dzieje
się to od czterdziestu do siedemdziesięciu razy. Efekt jest taki, że
przypominanie z pamięci nie jest czymś, co musisz sobie sam zorganizować na
końcu rozdziału \dash{} jest jedynym sposobem, w jaki da się przewracać strony.

Konstrukcja stopniuje też własne wsparcie: przykład rozwiązany, potem przykład
z lukami, potem \emph{jeden do wykonania dla ciebie}, potem zadania
samodzielne. I celowo zastawia pułapki. Tam, gdzie temat ma typowe błędne
przekonanie, tekst wprowadzi cię w nie, pozwoli ci się zadeklarować przy złej
odpowiedzi, a potem stwierdzi wprost, że jest błędna, i pokaże dlaczego. To
jest niewygodne i jest najcenniejszą rzeczą w tym formacie, bo błędne
przekonanie wywołane i skorygowane znika, a takie, którego jedynie uniknięto,
czeka.

\section*{Jak książka jest zbudowana}

Dziewięć części, czterdzieści sześć programów.

\textbf{Część~I \dash{} Podstawy} (F1--F13) nie zakłada niczego i jest sortowana
Quizami. Jest nachylona ku temu, co się później opłaci, a nie ku szkolnemu
programowi: logarytmy dostają cały program, bo arytmetyka w skali
logarytmicznej jest nośna trzy części dalej, a reguła łańcuchowa dostaje
najdłuższy program w części, bo propagacja wsteczna \emph{jest} regułą
łańcuchową i niczym więcej.

\textbf{Część~II \dash{} Liczba, precyzja i koszt} (P1--P3) wyprzedza algebrę
liniową, co jest odejściem od każdego programu nauczania matematyki i jest
zamierzone. Wszystko po niej to arytmetyka wykonywana przez skończoną maszynę
przy ograniczonym budżecie. Książka, która czeka czterysta stron, zanim
przyzna, że maszyna nie potrafi przedstawić liczby $\num{0.1}$, przez czterysta
stron uczyła fikcji.

\textbf{Część~III \dash{} Algebra liniowa} (P4--P10) jest największa, bo to
dziedzina, z której ten czytelnik korzysta codziennie i którą rozumie
najsłabiej.

\textbf{Część~IV \dash{} Struktury dyskretne i argumentacja} (P11--P13) obejmuje
zliczanie, grafy oraz czytanie twierdzenia bez wierzenia w więcej, niż ono
mówi. Stoi przed rachunkiem różniczkowym, żeby \emph{DAG} i \emph{dla każdego
$\varepsilon$} były zdefiniowane wtedy, gdy są pierwszy raz potrzebne.

\textbf{Część~V \dash{} Rachunek różniczkowy i różniczkowanie automatyczne}
(P14--P17) prowadzi regułę łańcuchową aż do tego, co robi \code{loss.backward()}.

\textbf{Część~VI \dash{} Optymalizacja} (P18--P21) w pierwszym tomie Strouda nie
występuje w ogóle.

\textbf{Część~VII \dash{} Prawdopodobieństwo i statystyka} (P22--P27) wychodzi poza
rozkład normalny, na którym Stroud się zatrzymuje, ku wnioskowaniu i metodzie
bayesowskiej \dash{} bo prawdziwym statystycznym zadaniem tego czytelnika jest
rozstrzygnięcie, czy zmierzona różnica jest prawdziwa, a tego nie obejmuje
żaden podręcznik matematyki inżynierskiej.

\textbf{Część~VIII \dash{} Teoria informacji} (P28--P30) to miejsce, z którego
bierze się funkcja straty.

\textbf{Część~IX \dash{} Złożenie} (P31--P33) poświęca całą książkę jednej
architekturze, jednemu treningowi i jednemu pomiarowi, i nie wprowadza nowej
matematyki.

\section*{Dwie zasady, których ta książka pilnuje}

\begin{note}
\textbf{Każda liczba jest policzona, nie zapamiętana.} Każda wartość liczbowa
tutaj wydrukowana pochodzi ze skryptu w katalogu \code{code/} i jest wciągana
do tekstu mechanicznie. Książka nie zawiera cyfr; zawiera odwołanie do nich.
Wartość, której skrypty przestaną produkować, drukuje widoczny znacznik i
przerywa budowanie. Istnieje to dlatego, że błędna cyfra w książce
matematycznej nie jest zepsutym przykładem \dash{} jest czymś, co czytelnik ze
sobą zabierze.
\end{note}

\begin{note}
\textbf{To, czego nie zmierzono, jest oznaczone jako osąd.} W tych programach
opisano dziesięć eksperymentów. Tam, gdzie któregoś nie przeprowadzono, teza,
którą by wspierał, mówi o tym wprost, a tabela, która miałaby zawierać jego
liczby, pozostaje pusta. Puste tabele nie są przeoczeniem \dash{} są mechanizmem.
\end{note}

\section*{Dwa języki, jedno źródło}

Ta książka istnieje po polsku i po angielsku i jest budowana z jednego
repozytorium. Oba wydania dzielą strukturę, rysunki, numerację zadań i
odpowiedzi; różni je wyłącznie tekst. Kontrola przy każdym budowaniu dowodzi,
że oba mają te same programy, te same liczby ramek i te same liczby, więc
poprawka wprowadzona w jednym wydaniu nie może po cichu ominąć drugiego.

Tam, gdzie polski i angielski uzus matematyczny naprawdę się różnią \dash{}
$\mathrm{tg}$ wobec $\mathrm{tan}$, przecinek dziesiętny, nawiasy ostre polskiej tradycji dla
przedziału domkniętego, $D^2(X)$ wobec $\Var(X)$ \dash{} decyzja jest zapisana w
Dodatku~B, a nie zostawiona pamięci tłumacza, i jest wykonywana przez skład, a
nie stosowana ręcznie.
